\section{Native time and coherent quantum transport}
\label{app:m1-quantum-transport}

The local-clock proposal gives classical records a complete influence cone.
A coherent extension must also charge the physical time of intermediate
motion. We characterize a finite register-flight and onsite-scattering
class, construct its optimal isotropic Weyl limit, retain every low-energy
species, and identify the stronger coherent structure needed for a shared
matter and record light speed. This is a theorem about a declared operation
class, not an admission certificate for a complete A1--A3 source.

\paragraph{The native flight class.}
Let a flight of duration $t$ have orthogonal internal projectors $P_s$
summing to the identity and displacements $b_s$ with $\|b_s\|\leq ct$.
Its Fourier action is $F(k)=\sum_s e^{-ik\cdot b_s}P_s$.
The modes move along their straight segments during the flight. Onsite
unitaries and waits may be inserted; all sequential stage durations add.
Internal dimension is arbitrary and finite, and every intermediate event
is retained. Flights move coherent modes without copying unknown states.
Their exterior-power lifts act on fermionic Fock space. Arbitrary admitted
record-flight directions attain the round cone of speed $c$; a fixed
direction menu has its own polyhedral reachability cone.

\paragraph{Sharp velocity theorem.}
For a word $U(k)$ of total duration $T>0$, let the isometry $V$ select a
two-dimensional degenerate eigenband at $k=0$. Its first-order velocity
matrices are
\[
 A_i=\frac{i}{T}V^*U(0)^*(\partial_iU)(0)V
     =\sum_a v_{a,i}E_a,\qquad
 E_a\geq0,\quad \sum_aE_a=I_2,\quad \|v_a\|\leq c.
\]
Indeed differentiating the word produces positive compressions of each
flight projector, weighted by that stage's fraction of the actual duration.
Waits and onsite time contribute zero velocities. Ancillary states and
quantum interference remain in the compressed effects.
Write $E_a=w_a(I+r_a\cdot\sigma)$, where $w_a\geq0$,
$\|r_a\|\leq1$ and $\sum_a w_a=1$. For
$A_i=u_iI+\sum_jM_{ij}\sigma_j$,
\[
 M=\sum_a w_av_ar_a^{\mathsf T},\qquad
 \|M\|_*\leq\sum_a w_a\|v_a\|\|r_a\|\leq c.
 \tag{QT1}
\]
Here $\|\cdot\|_*$ is the sum of singular values. Conversely, any real
matrix $M=\sum_j s_j u_jr_j^{\mathsf T}$ with $\sum_j s_j\leq c$
is attained without tilt: perform conditional flights with velocities
$\pm cu_j$ and projectors $(I\pm r_j\cdot\sigma)/2$ for fractions
$s_j/c$ of the period, then wait. Thus (QT1) is the complete untilted
first-order attainable set. General axes give finite addressed experiments;
periodic lattice realizations require commensurate directions or controlled
approximations. No characterization of the joint tilted pair $(u,M)$ is
asserted.

In particular an isotropic matrix $M=vR$, $R$ orthogonal, requires
$3|v|\leq c$. A normalized trace witness gives the same bound for a
three-component Clifford band: $(n\cdot\alpha)^2=\|n\|^2I$ bounds
each positive-effect contribution. More coin states, directions, local gates
or grouped periods cannot remove this factor within the stated grammar and
regular linear band limit. This excludes an isotropic Weyl or Dirac sector
with the complete record-cone speed, rather than merely failing to find one.
It is not a theorem about every quantum cellular automaton.

\paragraph{Saturation derives the minimal geometry.}
For $M=(c/3)I$, equality in the trace bound requires every nonzero effect
to have $\|v_a\|=c$, $\|r_a\|=1$ and $v_a=cr_a$. Consequently
\[
 \sum_a w_an_a=0,\qquad
 \sum_a w_an_an_a^{\mathsf T}=I/3,\qquad \|n_a\|=1.
 \tag{QT1a}
\]
Conversely these weighted spherical two-design identities attain the bound.
They are derived from optimality. At least four nonzero outcomes are needed:
the three independent coordinate columns are orthogonal to the weight
column. With exactly four, the matrix with rows
$\sqrt{w_a}(1,\sqrt3 n_a)$ has orthonormal columns and therefore rows.
The row norms force $w_a=1/4$ and orthogonality forces
$n_a\cdot n_b=-1/3$ for $a\ne b$. Thus the regular equally weighted
tetrahedron is the unique minimal optimizer, up to rotation, reflection and relabeling.
This minimum counts distinct velocity outcomes across the whole word.
A single flight therefore needs at least four internal channels, since
its distinct nonzero projectors are orthogonal. A multi-stage word can
reuse fewer channels: the three-axis example below uses two, with six
velocity outcomes. None of these counts is an OPH observer population or
a choice of central source ports.

\paragraph{An explicit minimal four-channel process.}
Use displacements $a s$ with
\[
 s\in\{(1,1,1),(1,-1,-1),(-1,1,-1),(-1,-1,1)\}
\]
and the onsite coin
\[
 C=\frac{1}{\sqrt3}
 \begin{pmatrix}
 0&1&1&1\\1&0&i&-i\\1&-i&0&i\\1&i&-i&0
 \end{pmatrix},\qquad C^*=C,\quad C^2=I.
\]
Apply $C$, then one simultaneous flight of the four channels at speed $c$.
The duration is $T_a=\sqrt3 a/c$ and
$U_{\rm tet}(\theta)=\operatorname{diag}(e^{-i\theta\cdot s})C$.
The projectors $P_\pm=(I\pm C)/2$ have rank two. For
$B_i=\operatorname{diag}(s_i)$, the matrices
$X_i^\pm=\sqrt3 P_\pm B_iP_\pm$ obey Pauli multiplication on their
bands, with opposite orientations. The compressed flight effects are
$(P_\pm+\sum_i s_iX_i^\pm/\sqrt3)/4$.
Hence the physical band speed is $(a/\sqrt3)/T_a=c/3$.
The native clock, coherent mixing and minimal direction geometry are all
specified in the same process.
It also transmits an exact native record at speed $c$: in an empty-register
background, leave the source empty for zero and prepare the local
one-particle state $C|s\rangle$ for one. The coin gives
$C^2|s\rangle=|s\rangle$, and the flight reaches $as$ in time
$\sqrt3a/c$. The receiver occupation effect distinguishes them with
probabilities zero and one. This attained response belongs to the same
walk. These inputs do not replace the fixed filled reference used below.
The experiment uses the complete finite channel family, without a uniform
continuum quasienergy bound on the preparation. Restricting interventions
to a low-energy sector would change the complete-clock requirement.

Its entire four-channel characteristic polynomial is
\[
 \det(\lambda I-U_{\rm tet})=
 \lambda^4-\frac{2}{3}(\cos2x+\cos2y+\cos2z)\lambda^2+1.
\]
Writing $\cos2\epsilon=(\cos2x+\cos2y+\cos2z)/3$ with
$0\leq\epsilon\leq\pi/2$, the eigenvalues are
$\pm e^{\pm i\epsilon}$. Exactly eight zero-quasienergy cones occur at
$\{0,\pi\}^3$, with parity selecting four cones of each chirality.
The same momenta contain a $\pi$ band; the eight points
$\{\pi/2,3\pi/2\}^3$ are degeneracies at $\pm i$, not extra zeros.
With the fixed logarithm convention below, both energies
$\epsilon/T_a$ and $(\pi-\epsilon)/T_a$ occur twice at every momentum.
The identity $\sin^2\epsilon=\sum_i\sin^2\theta_i/3$ gives
$\epsilon\geq2d(\theta,\{0,\pi\}^3)/(\pi\sqrt3)$.
It proves the same complete thermal limit (QT5); the second band is bounded
below by $\pi/(2T_a)$ and its full thermal contribution vanishes in the
refinement limit. No finite channel is discarded. Independent checks include
all four-channel momentum matrices, both bands, the full sixteen-dimensional
onsite Fock coin and all charged trajectories.

There is also a full four-channel dynamical bound, without selecting a
momentum-dependent band. Put $B=\operatorname{diag}(s\cdot k)$ and
$X_i=X_i^++X_i^-=\sqrt3(B_i+CB_iC)/2$.
Exactly $U_{\rm tet}(ak)^2=e^{-iaB}e^{-iaCBC}$; the product error is
at most $a^2\|[B,CBC]\|/2\leq3a^2\|k\|^2$.
For $j=\lfloor t/T_a\rfloor$, $\|k\|\leq K$ and $0\leq t\leq T$,
\[
 \|U_{\rm tet}(ak)^j-C^j e^{-it(c/3)k\cdot X}\|
 \leq \frac{\sqrt3 cT}{2}aK^2+\frac4{\sqrt3}aK.
\]
The final unpaired tick costs at most $\sqrt3aK$ and time rounding at
most $aK/\sqrt3$. At other nodes replace $C$ by $(-1)^{\rm parity}C$.
The actual zero band has unit carrier, while the other band keeps its
alternating phase. Omitting this carrier makes an order-one error at odd
ticks even at $k=0$. Blocking two ticks cannot turn the retained $\pi$
quasienergy sector into another physical zero-energy species.

\paragraph{A bound for arbitrary finite stencils.}
The flight factorization is not needed for a second obstruction. Let
$U(k)=\sum_{b\in S}e^{-ik\cdot b}K_b$ be any translation-invariant
finite-range unitary with finite internal dimension and period $T$.
For $G(n)=iU(k)^*\partial_nU(k)$, $l=\min_b n\cdot b$ and
$h=\max_b n\cdot b$,
\[
 lI\leq G(n)\leq hI.
 \tag{QT1b}
\]
To prove this, take a lattice-rational direction and rescale to integer
exponents. The polynomial
$V(z)=\sum_bz^{n\cdot b-l}e^{-ik\cdot b}K_b$ is unitary on the circle
and contractive in the disk by the maximum principle. The radial derivative
of $\|V(r)\psi\|^2$ at $r=1$ is nonnegative, giving the first inequality.
The reversed polynomial gives the second. Rational-direction density gives
the result for every $n$. Neither positivity nor orthogonality of the
coefficient matrices is assumed.

Compression to a degenerate two-state band implies
$u+M\overline B_1\subseteq\operatorname{conv}(S/T)$.
An untilted isotropic speed $v$ therefore requires a ball of radius $|v|T$
inside that polytope. Exact native speed $c$ requires $\|b\|\leq cT$.
A finite polytope inside the $cT$ ball cannot contain the same round ball,
so finite stencils cannot attain an exactly isotropic speed $c$.
The caps $\{n\in S^2:n\cdot b\geq |v|T\}$ must cover the
unit sphere when $|v|>0$. Each has area at most $2\pi(1-|v|/c)$. If $N$ counts nonzero
displacement terms, sphere area and subadditivity give
\[
 N(1-|v|/c)\geq2.
 \tag{QT1c}
\]
Thus fixed $N$ excludes even a refinement limit with $|v|/c\to1$;
accuracy $|v|/c\geq1-\eta$ requires $N\geq2/\eta$.
Increasing internal dimension alone cannot evade it. The bound also
applies to approximately isotropic refinements: use
$s_{\min}(M)-\|u\|$ in place of $|v|$ when it is strictly positive,
since that centered ball lies in the velocity ellipsoid. A nonpositive
value supplies no positive-radius cap bound. For example,
$U(k)=e^{-ik_x}I_2$ with $T=c=1$ has $N=1$, $u=(1,0,0)$ and $M=0$;
its velocity set does not contain the origin. Clamping the radius to zero
would falsely demand $1\geq2$. The stationary identity has $N=0$ and
is also outside this positive-radius argument. Exact scope controls retain
both processes. For $M\to cR$, $u\to0$, the radius is eventually positive
and tends to $c$, so $N\to\infty$ follows without exact isotropy
at any cutoff. The tetrahedral
process saturates its support-polytope bound as well as the flight bound.
The cube has inradius $a$ and circumradius $\sqrt3a$, so even an
endpoint-only clock for the next example gives at most $c/\sqrt3$;
its actual sequential flights give $c/3$. A growing direction set is
necessary, but does not prove existence of a luminal unitary refinement.
The stronger flight bound excludes such a limit for every positively timed
register-flight lowering, even with growing direction sets.

\paragraph{An optimal native process.}
On the unit three-torus take $q\in4\mathbb N$, $a=1/q$, and two fermionic
modes at each site. Apply conditional axis flights of length $a$, with
Pauli projectors along $x,y,z$ in that order. Each costs $a/c$, so
\[
 T_a=3a/c,\qquad
 U(\theta)=e^{-i\theta_z\sigma_z}e^{-i\theta_y\sigma_y}
               e^{-i\theta_x\sigma_x}.
 \tag{QT2}
\]
Translations and onsite basis changes realize these exact unitaries.
For physical momentum $k=\theta/a$, Taylor remainders and unitary
telescoping give, for $\|k\|\leq K$ and $0\leq t\leq T$,
\[
 \left\|U(ak)^{\lfloor t/T_a\rfloor}
       -e^{-it(c/3)k\cdot\sigma}\right\|
 \leq \frac{2cT}{3}aK^2+aK.
 \tag{QT3}
\]
The product remainder is at most $a^2(\sum_i|k_i|)^2/2$, the single
exponential remainder at most $a^2\|k\|^2/2$, and the final incomplete
period costs at most $aK$. The construction attains the speed bound.
Assigning the whole word duration $a/c$ would omit two actual flights.
Onsite overhead can only lower the speed. An overhead fraction tending to
zero preserves (QT3) in the limit; uniformly bounded control strength is
not asserted to realize arbitrarily fast fixed-angle basis changes.

\paragraph{Complete spectrum and fixed reference.}
The eigenvalues are $e^{\pm i\epsilon}$, $0\leq\epsilon\leq\pi$, with
\[
 \cos\epsilon=\cos x\cos y\cos z+\sin x\sin y\sin z.
 \tag{QT4}
\]
There are exactly eight zeros: four points in $\{0,\pi\}^3$ with even
$\pi$ parity, and four in $\{\pi/2,3\pi/2\}^3$ with even $3\pi/2$
parity. Cauchy--Schwarz gives the bound
\[
 |\cos\epsilon|\leq\sqrt{\cos^2x\cos^2y+\sin^2x\sin^2y}\leq1.
\]
Equality forces
$\cos x\sin y=\sin x\cos y=0$ and the stated $z$ alignment.
The opposite parities give eight $\pi$ points. At the first four zeros
the linear Pauli coefficient matrix is $I$; at the other four it is
$\operatorname{diag}(1,-1,1)$. All eight Weyl cones propagate at $c/3$,
four with each chirality. None is removed.

Define the finite reference using the principal Hermitian logarithm of the
actual $U$, divided by $T_a$, with a fixed branch at $-I$. Fill strictly
negative levels and leave zero levels empty. The Fock lift preserves this
Slater vacuum exactly. Normal ordering gives two excitation choices per
momentum of energy $E_a=\epsilon/T_a$. This is a declared Floquet
quasienergy; the generally nonlocal logarithm is not a derived autonomous
Hamiltonian or an OPH physical energy calibration.

For every $\beta>0$, the complete finite partition function satisfies
\[
 \log Z_a=2\sum_\theta\log(1+e^{-\beta E_a(\theta)})
 \longrightarrow
 16\sum_{n\in\mathbb Z^3}
       \log(1+e^{-\beta(2\pi c/3)\|n\|}).
 \tag{QT5}
\]
The normal-ordered energy and Gibbs entropy converge as well. An explicit
global bound is $\epsilon(\theta)\geq d(\theta,Z)/(8\sqrt3)$.
For $\epsilon\geq\pi/8$, the torus diameter proves it. Otherwise
$\cos(2x)\cos(2y)\geq\cos(2\epsilon)>0$ puts $x,y$ within
$\epsilon$ of the same zero family. Writing
$R e^{i\phi}=\cos x\cos y+i\sin x\sin y$, one obtains
$|z-\phi|\leq\epsilon$ and an additional angular deviation at most
$\arctan(\tan^2\epsilon)\leq\epsilon$, hence
$d(\theta,Z)\leq\sqrt6\epsilon$.
Assigning each momentum to a nearest zero gives the uniform summable bound
$E_a\geq\pi c\|n\|/(12\sqrt3)$. Dominated summation and the eight
linearizations prove (QT5), including all zero-level degeneracy.
Fixed finite-excitation dynamics follow from exterior-product telescoping
and (QT3). Finite smeared CAR covariances use projector convergence away
from the explicitly treated zero levels. No global norm identification of
infinite filled seas or interacting-field limit is claimed.

\paragraph{Why a finite Hamiltonian does not repair exact delay.}
Let $H$ act on a finite tensor product of full matrix site factors at
distinct fixed positions. If every local pair satisfies
\[
 [e^{itH}A_xe^{-itH},B_y]=0\quad\text{for }0<t<\|x-y\|/c,
 \tag{QT6}
\]
analyticity makes that commutator zero for all $t$. The evolving algebra
at $x$ therefore stays in that factor. Its derivative is an inner onsite
derivation; subtracting all onsite implementers leaves a scalar. Thus
$H$ has no intersite interaction. Finite quadratic canonical propagators
have the same analytic obstruction. The hypothesis is exact zero influence
for all admitted local preparations and effects, not a detection threshold.

For two qubits and interaction coefficients $h_{ab}$ in the Pauli basis,
the complete derivative witness is
\[
 \sum_{i,j=1}^3
 \|[[H,\sigma_i\otimes I],I\otimes\sigma_j]\|_{\rm HS}^2
 =256\sum_{a,b=1}^3h_{ab}^2.
\]
For $H=\sigma_x\otimes\sigma_x$, first-qubit preparation $+y$ and
remote preparations $\pm x$ give receiver-effect probabilities
$(1\pm\sin2t)/2$. Thus an actual local preparation produces an
arbitrarily early signal. Moving modes evade fixed algebra ownership;
sampled automata lack the assumed continuous interpolation; continuum
fields use infinite local algebras and singular generators.

\paragraph{The coherent alternative.}
If the response vector space acts on an operationally supplied fundamental
SU(2) module, covariance forces a real-linear Hermitian velocity map to be
$A(v)=\gamma v\cdot\sigma$: Hermitian matrices split into the scalar and
adjoint representations. A covariant onsite term may add $\mu I$;
covariance does not determine that energy offset. Choosing zero onsite
offset, handedness and scale gives
$i\partial_t\psi=-ic\sigma\cdot\nabla\psi$. Its generator squares
to $-c^2\Delta$, and its conserved current obeys
$\rho=\psi^*\psi$, $j=c\psi^*\sigma\psi$, $\|j\|\leq c\rho$.
The exterior-ball flux identity proves exact propagation in radius $ct$,
with dispersion $\pm c\|k\|$. A parity-paired Dirac module permits an
anticommuting onsite mass without changing this cone. Standard CAR
quantization supplies locality of even observables and positive excitation
energies. These noncommuting luminal velocities cannot have the positive
flight resolution above, which would imply $3c\leq c$.

This standard continuum witness identifies a precise structure to realize;
it does not supply the fundamental module's source meaning, coherent
spatial generator, physical clock or refining local instruments. Internal
gauge SU(2) is not thereby identified with spin. Neither a new core axiom
nor a finite twelve-port realization is asserted.

\paragraph{Evidence and relation to prior mathematics.}
The package at \ophid{code/m1_quantum_transport} independently reconstructs
all six flight orders, exact momentum matrices, complete thermal catalogs,
all sectors and intermediate stages of finite fermionic gate programs,
positive velocity resolutions, Pauli delay probes and coherent currents.
Hostile controls reject uncharged time, hidden species, changed phases and
regenerated signless permutations. Thirteen Lean lemmas and a transitive axiom
audit certify finite algebra reductions; the operator, SVD and asymptotic
arguments above remain analytic proofs.
Quantum-walk relativistic limits are standard; see Mlodinow and Brun,
\emph{Physical Review A} 97, 042131 (2018),
\href{https://arxiv.org/abs/1802.03910}{arXiv:1802.03910}, and Bisio,
D'Ariano and Perinotti,
\href{https://arxiv.org/abs/1608.02004}{arXiv:1608.02004}.
The tetrahedral qubit measurement is a standard SIC construction; see
Renes, Blume-Kohout, Scott and Caves,
\href{https://arxiv.org/abs/quant-ph/0310075}{Journal of Mathematical Physics 45, 2171 (2004)}.
Tetrahedral Dirac walks are also studied by Nzongani and collaborators,
\href{https://arxiv.org/abs/2404.09840}{arXiv:2404.09840}.
The strict-causal versus Hamiltonian distinction is discussed by Ranard,
Walter and Witteveen,
\href{https://doi.org/10.1007/s00023-022-01193-x}{Annales Henri Poincar\'e (2022)}.
The contribution here is the sharp native-time characterization, its
optimal complete-spectrum construction and their implication for the
proposed common record and matter clock.
